\documentclass[../../../include/open-logic-section]{subfiles}

\begin{document}

\olfileid[fr]{sth}{story}{urelements}
\olsection{Des uréléments ou non ?}

Dans les prochains chapitres, nous chercherons à dégager des axiomes
de la conception cumulative et itérative des ensembles. Mais, avant
d'aller plus loin, il nous faut préciser le rôle des \emph{uréléments}.

Le schéma de la \olref[approach]{sec} ne nous permettait de former de
nouveaux ensembles qu'à partir d'\emph{ensembles déjà formés}. Nous
pourrions toutefois vouloir admettre que certains \emph{objets qui ne
sont pas des ensembles} — vaches, cochons, grains de sable ou autres —
soient des !!{element}s d'ensembles. Dans ce cas, nous partirions
d'objets fondamentaux, les \emph{uréléments}, puis nous dirions qu'à
chaque étape $S$ nous formons «~tous les ensembles possibles~» dont
les membres sont des uréléments ou des ensembles formés aux étapes
antérieures à $S$ (dans n'importe quelle combinaison). Le schéma
obtenu ressemblerait plutôt à ceci :
\begin{center}
	\begin{tikzpicture}[scale=0.6]
	\tikzset{cut_here/.style={densely dotted}}
	\draw (-4,6) -- (-1, 0) -- (1,0) -- (4,6);
	\draw[cut_here] (-5,8)--(-4,6);
	\draw[cut_here] (5,8)--(4,6);
	\draw(-1.5, 1)--(1.5, 1);
	\draw(-2, 2)--(2, 2);
	\draw(-2.5, 3)--(2.5,3);
	\draw(-3, 4)--(3, 4);
	\draw(-3.5, 5)--(3.5, 5);
	\draw(-4, 6)--(4, 6);
	\node[label] at (2, 0) {\small 0};
	\node[label] at (2.5, 1) {\small 1};
	\node[label] at (3, 2) {\small 2};
	\node[label] at (3.5, 3) {\small 3};
	\node[label] at (4, 4) {\small 4};
	\node[label] at (4.5, 5) {\small 5};
	\node[label] at (5, 6) {\small 6};
	\end{tikzpicture}
\end{center}
Nous devons donc trancher : \emph{faut-il admettre des uréléments ?}

Du point de vue philosophique, leur inclusion se défend. La raison
principale est de rendre notre théorie des ensembles \emph{applicable}.
Pour illustrer cette idée, rappelons que, dans le
\olref[sfr][siz][]{chap}, nous avons dit que deux ensembles $A$ et~$B$
ont la même taille, c'est-à-dire que $\cardeq{A}{B}$, si et seulement
s'il existe une bijection entre eux. Or, si les vaches d'un pré
forment un ensemble et les cochons d'une porcherie en forment un
autre, nous pouvons traiter en termes ensemblistes l'affirmation
«~il y a autant de vaches que de cochons~». Si, en revanche, nous
excluons les uréléments, de sorte que les vaches et les cochons ne
forment \emph{pas} d'ensembles, ce traitement ensembliste n'est plus
disponible. Nous ne pouvons alors appliquer directement la théorie
des ensembles qu'aux ensembles eux-mêmes. (Pour d'autres arguments
en faveur des uréléments, voir \citealt[pp.~vi, 24,
50--1]{Potter2004}.)

En mathématiques, cependant, il est assez rare d'admettre des
uréléments. C'est en partie parce qu'il est \emph{un tout petit peu}
plus facile de formuler la théorie des ensembles sans eux. On trouve
néanmoins parfois des raisons plus intéressantes de les exclure :
\begin{quote}
	Si l'on estime que la théorie des ensembles fonde les
	mathématiques, nous devrions pouvoir saisir toutes les
	mathématiques en ne parlant que d'ensembles ; nos variables ne
	devraient donc pas parcourir des objets tels que les vaches et les
	cochons.
	%Mais si $C$ est une vache, $\{C\}$ est un ensemble sans être un objet mathématique légitime.
	\citep[p.~8]{Kunen1980}
\end{quote}
Ainsi, le souci de l'applicabilité inciterait à \emph{admettre} des
uréléments ; une visée fondatrice réductrice, qui ramène les
mathématiques à la théorie pure des ensembles, pourrait inciter à
les \emph{exclure}. Une légère paresse conduit aussi à les exclure.

Nous suivrons la voie la plus paresseuse. Il y a toutefois à cela une
raison pédagogique : nous voulons présenter les rudiments de la
théorie des ensembles nécessaires pour aborder sa philosophie. Or
la majeure partie de la littérature philosophique sur ce sujet
traite de théories des ensembles formulées \emph{sans} uréléments.
Ce livre sera donc peut-être plus utile s'il reste assez proche de
cette littérature.

\end{document}
